
\section{模型优缺点评价}

% 所谓的模型优缺点评价往往并不局限于模型本身,在整个建模过程中所表露出的优缺点均可在最后进行陈述,一般撰写模型优缺点的基本原则是优点说充分,缺点不回避。

\subsection{模型的优点}

常见的优点表述形式如下:

\begin{itemize}[itemindent=2em]
  \item \textbf{优点 1：模型或思路设计的简洁实用,效率高。}【结合本文说明】~\cite{ref1}。
  \item \textbf{优点 2：本文建立的模型具有很强的创新性。}【结合本文说明】。
  \item \textbf{优点 3：模型的计算结果准确,精度高。}【结合本文说明】。
  \item \textbf{优点 4：模型考虑的系统全面,有很强的实用性。}【结合本文说明】。
  \item \textbf{优点 5：对模型进行了各类检验、稳定性高。}【结合本文说明】。
  \item \textbf{优点 6：模型本身具有的优点。}【结合本文说明】。
\end{itemize}

\subsection{模型的缺点}

常见的缺点表述形式如下:

\begin{itemize}[itemindent=2em]
  \item \textbf{缺点 1：}受【XX 因素】限制,未考虑【XX 情况】,影响精度。
  \item \textbf{缺点 2：}本文考虑的因素较为理想,降低了模型的普适性和推广能力。
  \item \textbf{缺点 3：}由于系统考虑了【XXX】等因素,导致模型较为复杂,计算时间长,效率低。
  \item \textbf{缺点 4：}模型本身具有的缺点。
\end{itemize}

\subsection{模型的改进}

% 模型的改进一般是针对模型的缺点而言的,主要是提出一些改进的思路即可。

针对上述不足,本文可从以下方面进行改进:【改进方向 1】,通过【具体方法】以解决【具体问题】;【改进方向 2】,通过【具体方法】以提升【具体指标】;【改进方向 3】,通过【具体方法】以扩展【应用场景】。后续研究将围绕上述方向进一步深入,使模型在更广泛的实际场景中保持高精度与强鲁棒性。
